%%%%%%%%%%%%%%%%%%%%%%%%%%%%%%%%%%%%%%%%%%%%%%%%%%%%%%%%%%%%%%%%%%%%%%%%%%%
%% Glossary entries (suggestion: append "g" to the label).
%%%%%%%%%%%%%%%%%%%%%%%%%%%%%%%%%%%%%%%%%%%%%%%%%%%%%%%%%%%%%%%%%%%%%%%%%%%

\newglossaryentry{GULg}{
    name={Linux Users Group},
    description={A student association interested in computers and computing,
                 sharing the common goal of using and promoting free software}
}

\newglossaryentry{UC3Mg}{
    name={Universidad Carlos III de Madrid},
    description={Spanish public university located in the Region of Madrid}
}

\newglossaryentry{swlibre}{
    name={Free software},
    description={Software that respects users' freedom to run, copy, study,
                 modify and redistribute it}
}


%%%%%%%%%%%%%%%%%%%%%%%%%%%%%%%%%%%%%%%%%%%%%%%%%%%%%%%%%%%%%%%%%%%%%%%%%%%
%% Acronym entries (suggestion: append "a" to the label).
%%  An acronym may have a glossary entry (GUL, UC3M) via see=; or not (OS),
%%  to illustrate both cases.
%%%%%%%%%%%%%%%%%%%%%%%%%%%%%%%%%%%%%%%%%%%%%%%%%%%%%%%%%%%%%%%%%%%%%%%%%%%

\newglossaryentry{GULa}{
    type = \acronymtype,
    name = {GUL},
    description = {Linux Users Group},
    see=[Glossary:]{GULg}
}

\newglossaryentry{OSa}{
    type = \acronymtype,
    name = {OS},
    description = {Operating System}
}

\newglossaryentry{UC3Ma}{
    type = \acronymtype,
    name = {UC3M},
    description = {Universidad Carlos III de Madrid},
    see=[Glossary:]{UC3Mg}
}
